\chapter{Conclusion and Next Activities}

\section{Conclusion}
This work developed and evaluated a hybrid framework for mission-driven ROS~2 configuration. The framework combines static repository analysis, an explicit system model, bounded language-model reasoning, deterministic capability realization, validated parameter modification, and template-based artifact generation. Its central design decision is to use the language model only where semantic interpretation is required, while keeping repository facts, dependency resolution, admissible configuration values, file structure, and rendering under deterministic control. The implemented prototype analyzes ROS~2 Python nodes, launch descriptions, and parameter files without executing them. It integrates the extracted facts into a deployment model containing component activation, parameter provenance, interfaces, remappings, and known communication relationships.

Natural-language reasoning is divided into three specialized calls executed in sequence: capability reasoning selects an implementation for each supported capability, parameter selection chooses identifiers only from a capability-conditioned, top-ten semantic retrieval context, and parameter-value reasoning receives only the validated selection and generates typed replacement values. The model never writes YAML or launch code directly. Structured responses are checked against schemas, active-component restrictions, retrieved identifiers, types, allowed values, numerical bounds, and parameter relationships before deterministic templates can materialize a configuration.

The complete pipeline was evaluated on the train (150 missions) and test (100 missions) datasets described in Chapter~\ref{chap:dataset-results}, which span supported, unsupported, and ambiguous outcomes. Capability interpretation, evaluated in isolation, reached 90.0\%/92.0\% exact match. Parameter selection and value reasoning, evaluated in isolation given each mission's correct capability context, reached 76.5\%/86.7\% selection exact match and 87.3\%/92.5\% individual-value accuracy. Evaluated as a complete, live pipeline, in which each call receives only the outputs of the calls before it, joint exact match reached 65.6\%/66.0\%, and the unsafe execution rate was 26.7\%/38.0\%.

The principal contribution is therefore not unrestricted natural-language control of a robot, but a reproducible and auditable boundary between probabilistic interpretation and deterministic system modification. Repository-derived evidence limits what the model can select; capability realization determines what may be active; validation prevents malformed or ungrounded conclusions from reaching the renderer; and templates preserve syntactic correctness and unchanged baseline values. The results in Chapter~\ref{chap:dataset-results} show this boundary holding even where the underlying language-model predictions are wrong: every recorded failure is a wrong or missing decision that is rejected or scored incorrect, not a malformed or ungrounded change reaching the renderer.

Operational reliability has not yet been established. Every reported experiment is a single run of one quantized local model, so variance and confidence intervals are not yet available, and the train and test paraphrases remain in \texttt{pending} review. Structural validation of generated artifacts is also not equivalent to successful behavior in simulation or on a physical robot, which has not been attempted. The unsafe execution rate reported above is the clearest evidence that this system is not yet reliable: a substantial share of missions whose correct answer is \texttt{unsupported} or \texttt{needs\_clarification} are instead executed with a confidently produced configuration. The accuracy figures reported in this work should therefore be read as evidence for the proposed architecture on the datasets evaluated so far, not as a claim of safe or general natural-language autonomy.

\section{Next Activities}
\label{sec:next-activities}

The next activities should consolidate the current result before expanding the system:
\begin{enumerate}
    \item \textbf{Extend the outcome space.} Add independently annotated unsupported, contradictory, and genuinely ambiguous missions. This is necessary to measure unsafe acceptance, unsupported-request rejection, clarification precision and recall, and behavior when no valid configuration exists. Targeted development should focus on negation, shorthand, multi-clause requests, implicit units, coupled physical properties, and the distinction between selecting an implementation and tuning its parameters.

    \item \textbf{Extend the boundary breakdown of end-to-end errors.} Chapter~\ref{chap:dataset-results} already classifies joint-pipeline failures by which stage's output first diverges from the reference; extend this into a full root-cause analysis of each category, distinguishing runtime or validation errors from semantically valid but incorrect structured outputs, to determine whether the next improvement belongs in metadata, retrieval, prompting, validation, or the capability registry.

    \item \textbf{Run controlled ablations.} Hold the dataset, capability contract, model, and prompts fixed while independently varying deterministic evidence, LLM-derived aliases and user expressions, typed relationships, graph expansion, active-component filtering, candidate count, and validation-guided retry. Repeat the final experiment to estimate run-to-run variance and confidence intervals. Compare alternative local models under the same frozen prompts and contexts, reporting semantic accuracy, validation failures, latency, memory requirements, and token usage rather than accuracy alone.

    \item \textbf{Validate generated behavior in simulation.} Execute generated bundles in a simulator such as Gazebo, then follow with controlled experiments on the physical AntRobot where safe and feasible.
\end{enumerate}

These activities move the project from a repository-specific prototype toward a validated configuration assistant. The priority is to verify the results reported in this work under repeated and controlled conditions, and to demonstrate that accepted configuration plans produce the intended behavior in simulation and on the robot.
